\documentclass[sigplan,nonacm]{acmart}

\usepackage{booktabs}
\usepackage{listings}
\usepackage{xcolor}

\lstset{
  basicstyle=\ttfamily\small,
  keywordstyle=\bfseries,
  commentstyle=\itshape\color{gray},
  columns=fullflexible,
  keepspaces=true,
  literate={→}{$\rightarrow$}1 {←}{$\leftarrow$}1 {≡}{$\equiv$}1
           {≤}{$\leq$}1 {≥}{$\geq$}1,
}

\begin{document}

\title{A Semi-Persistent E-Graph with Native AC Canonization and Leapfrog AC Matching}

\author{R\'{e}mi Delmas}
\affiliation{%
  \institution{Amazon Web Services, Automated Reasoning Group}
  \country{USA}
}
\email{delmasrd@amazon.com}

\begin{abstract}\sloppy
We present SP-Egraph, an e-graph engine with three contributions.
First, pervasive semi-persistence built from a single vector
primitive with O(1) snapshots and O($k$) restore, enabling
equality saturation inside backtracking search algorithms.
Second, native support for associative~(A),
associative-commutative~(AC), and AC-idempotent~(ACI) operators
with an extension of leapfrog triejoin to multiset matching
using maximum partition semantics.
Third, proof logging via copy-on-first-re-canonization with
Euler-tour LCA for batch extraction.
Both semi-persistence and proof logging are compile-time opt-out
with zero residual overhead.
\end{abstract}

\maketitle

\section{Introduction}

E-graphs have become a central data structure in program optimization
and verification, driven by the practical success of
\textsc{egg}~\cite{egg} and \textsc{egglog}~\cite{egglog}.
Building on these foundations, and taking inspiration from the
AC(X) decision procedure in the Alt-Ergo SMT
solver~\cite{iguernlala}, SP-Egraph makes semi-persistence and
native algebraic theory support its primary design concerns.

Semi-persistence allows the e-graph to be embedded in backtracking
search algorithms: each decision point is a generation, and
rollback restores the e-graph to a prior state in O($k$) time,
where $k$ is the number of distinct cells modified since the
snapshot.  Native support for associative (A),
associative-commutative (AC), and AC-idempotent (ACI) operators
prevents the combinatorial blowup that occurs when these properties
are modeled as rewrite rules.

To our knowledge, leapfrog-style worst-case-optimal join algorithms
have not previously been adapted to AC matching.  SP-Egraph
integrates AC multiset decomposition directly into the relational
e-matching engine introduced by \textsc{egglog}~\cite{egglog,veldhuizen}.

\medskip

SP-Egraph's design rests on three contributions:

\emph{Pervasive semi-persistence.}  Every data structure is built
from a semi-persistent vector whose mark/restore
protocol logs only the first mutation to each cell per frame.
Inspired by Conchon and Filli\^{a}tre~\cite{conchon-filliatre}, but
with a first-write-wins invariant that bounds the diff by distinct
cells touched.  All higher-level containers (union-find, hash-cons
caches, node stores, registries) compose from this primitive and
inherit snapshot semantics automatically.

\emph{Leapfrog triejoin extended to A/AC/ACI.}  Egglog introduced
relational e-matching via leapfrog triejoin~\cite{veldhuizen} for
plain operators.  SP-Egraph extends the join framework with
subsequence matching for A, sub-multiset matching with maximum
partition semantics for AC, and subset matching for ACI, all
integrated into the same cost-based scheduling engine.

\emph{Proof logging with compile-time opt-out.}  A dual-parent-pointer
union-find with copy-on-first-re-canonization preserves original node
structure for proof reconstruction.  Euler-tour LCA enables
O($n$)-preprocessing, O(1)-per-query batch proof extraction.  A \texttt{const PROOFS: bool} generic
eliminates all proof machinery at compile time when proofs are not
needed.

\emph{Extensible literal model.}  The e-graph is parameterized over
the representation of literal values and their operations.  Users
provide parsing, interning, and evaluation functions for their own
types (machine integers, arbitrary-precision arithmetic, strings,
intervals, etc.)\ via a \texttt{LitModel} trait.  Literal values are
auto-lifted into the e-graph as dedicated lift nodes rather than
being baked into child references via union types as in
\textsc{egglog}~\cite{egglog}.  This keeps the node representation
uniform and allows new sorts to be added without modifying the
engine.

\section{Semi-Persistent E-Graph}

\subsection{The Diff-Log Protocol}

The entire e-graph state is built from semi-persistent vectors:
\begin{itemize}
\item \emph{Union-find:} two parent vectors (a path-compressed fast
  path and an uncompressed proof path) plus rank and justification
  vectors.
\item \emph{E-classes:} represented by a circular list buffer
  (arena-allocated) for e-class members with O(1) splicing on merge,
  and a sparse set (three vectors, swap-and-pop) for e-class
  representatives with O(1) add, remove, membership test, and direct
  enumeration of all current roots.
\item \emph{Use-lists:} singly-linked lists stored in a vector arena,
  tracking which parent e-nodes reference each e-class (the worklist
  for rebuild).
\item \emph{Node caches:} partitioned by node type (arity 0--3,
  commutative, A, AC, ACI, literal).  A global routing table maps
  global e-node ids to cache-local ids within each partition.  Each
  partition stores nodes in a semi-persistent vector alongside a
  transient \texttt{HashMap} index that is rebuilt from the vector
  after each restore.
\item \emph{Variadic children:} variable-arity nodes (A, AC, ACI)
  store headers in one vector and children in a separate pool vector.
\item \emph{Proof history:} when proof logging is enabled, original
  node contents before re-canonization are saved in append-only
  vectors (one per cache partition).
\item \emph{Registries and literal store:} semi-persistent maps
  (append-only vector + transient index) for sorts, operators, rules,
  and interned literal values.
\end{itemize}

All data structures are marked and restored together.  Marking
triggers a full rebuild first, so only fully congruence-closed states
are ever snapshotted.

Each vector records negative diffs: before overwriting a slot, the
old value is pushed onto a per-frame diff stack.  \texttt{mark()}
pushes a frame in O(1); \texttt{restore()} replays the diff log in
reverse in O($k$), where $k$ is the number of distinct cells
modified since the mark.  The first-write-wins invariant ensures each
slot is logged at most once per frame, bounding the diff size by the
size of the data structure --- but in practice much smaller,
typically proportional to the number of new e-nodes and merges in the
current generation.

For dense identifier types (31-bit ids), the capture flag is the MSB
of the stored representation, costing zero extra memory.

The implementation uses arena-style allocation with deterministic
truncation on restore, avoiding GC pauses.  This contrasts with
Conchon and Filli\^{a}tre's OCaml formulation~\cite{conchon-filliatre},
which relies on the GC to reclaim abandoned patch chains.

\subsection{Compile-Time Opt-Out}

A \texttt{const TRACK: bool} generic on the semi-persistent vector
controls all diff-log machinery.  When \texttt{false}, the compiler
eliminates capture flags, diff stacks, and frame bookkeeping
entirely, reducing the vector to a plain \texttt{Vec} with zero
residual overhead.  This allows the same generic code to be used in
both backtracking and non-backtracking contexts.

\subsection{Source-of-Truth vs.\ Derived State}

The e-graph separates source-of-truth containers (node stores,
union-find, literal store, registries) from derived containers
(hash-cons indices, sorted indexes for pattern matching).  Only
source-of-truth containers participate in the diff-log protocol.
Derived structures are rebuilt from source-of-truth after each
restore, avoiding diff-tracking overhead on high-churn index
structures.

\subsection{Stratified Negation}

Because \texttt{mark()} always snapshots a fully rebuilt state,
the generational structure directly supports stratified negation:
stratum $k$ saturates, then \texttt{mark()} rebuilds and snapshots;
stratum $k{+}1$ writes into a new generation while negative lookups
query the frozen snapshot.

\section{Native Algebraic Theories}

\subsection{Canonical Representations}

SP-Egraph handles A, C, AC, and ACI properties structurally through
canonical node representations, not rewrite rules.

\begin{table}[t]
\caption{Canonical forms by operator kind.}
\label{tab:canon}
\begin{tabular}{@{}llll@{}}
\toprule
Kind & Children & Canonical form & Merge \\
\midrule
Plain & $[c_0, \ldots, c_n]$ & Order preserved & In place \\
C     & $[c_0, c_1]$         & Sorted pair     & Re-sort \\
A     & $[c_0, \ldots, c_n]$ & Order preserved & In place \\
AC    & $[(id, m), \ldots]$  & Sorted by id    & Sum mults \\
ACI   & $[id, \ldots]$       & Sorted, unique  & Dedup \\
\bottomrule
\end{tabular}
\end{table}

AC nodes store sorted multisets of (id, multiplicity) pairs.  When
two entries acquire the same canonical id after a merge, their
multiplicities are summed.  ACI nodes store sorted sets; duplicates
are removed.  This prevents the combinatorial blowup that occurs
when commutativity and associativity are encoded as rewrite rules.

\subsection{Surface Syntax for Variadic Matching}

The surface language provides rest variables and splicing for
matching and constructing variadic (A/AC/ACI) terms.  Rest variables
bind the unmatched residual of a variadic operator.

\begin{lstlisting}
;; Sorts and operators with algebraic attributes
(sort Expr)
(function Add (Expr) Expr :assoc-comm)
(function Mul (Expr Expr) Expr)
(function Or (Expr) Expr :assoc-comm-idem)
(sort List)
(function Nil () List)
(function Concat (List) List :assoc)

;; AC: remove zero from sum (any multiplicity)
(rewrite (Add (Zero):k ..rest)
         (Add ..rest))

;; AC: distribute multiplication over sum
(rewrite (Mul a (Add ..rest))
         (Add ..{(Mul a x) for x in rest}))

;; ACI: remove false from disjunction
(rewrite (Or (False) ..rest)
         (Or ..rest))

;; A: remove empty list from concatenation
(rewrite (Concat ..pre (Nil) ..suf)
         (Concat ..pre ..suf))
\end{lstlisting}

Rest variables (\texttt{..rest}, \texttt{..pre}, \texttt{..suf})
bind the unmatched residual of a variadic operator.  For AC
operators, the multiplicity annotation \texttt{x:k} binds both the
element and its count.  Multiplicity constraints (\texttt{x:k>=2},
\texttt{x:k<5}) are compiled to intervals; when a multiplicity
variable appears multiple times in a pattern (non-linear), the
intervals are intersected and the match fails early if the
intersection is empty.


\section{Relational E-Matching}

\subsection{Index Construction}

Four sorted index families are rebuilt each saturation iteration:
\texttt{by\_op} (operator $\to$ nodes), \texttt{by\_repr}
(representative $\to$ nodes), \texttt{by\_child\_pos} ((class,
position) $\to$ parent nodes), and \texttt{by\_contains} (class
$\to$ variadic parents).  All are sorted vectors; binary search on
contiguous memory outperforms incrementally maintained trees.

\subsection{Leapfrog Triejoin}

Patterns compile to flat relational atoms joined by worst-case-optimal
leapfrog triejoin~\cite{veldhuizen} over the sorted indexes.  A
cost-based scheduler orders variables by estimated selectivity each
iteration.  The scheduler binds variables in any order (top-down,
bottom-up, or middle-out) because the flattened atom representation
decouples pattern structure from execution plan.

\subsection{Extending Leapfrog to AC Matching}

Standard leapfrog triejoin handles plain operators where each child
occupies a fixed position.  AC operators store sorted multisets of
(id, multiplicity) pairs, so the pattern matcher must decompose a
multiset into bound elements and a residual, rather than extracting
children by position.

The AC matching algorithm proceeds in two phases within the
backtracking DFS engine:

\paragraph{Phase 1: Join.}  The query plan first binds the AC
node's e-class representative via a standard leapfrog join over
\texttt{by\_op} and \texttt{by\_repr} indexes.  This identifies
which AC nodes to decompose.

\paragraph{Phase 2: Decompose.}  For each matched AC node, the
engine extracts the sorted multiset of (representative, multiplicity)
pairs and runs a recursive decomposition:

\begin{enumerate}
\item Initialize the \emph{residual pool}: the full multiset of
  (id, multiplicity) pairs after canonicalizing each child through
  \texttt{find}.

\item For each pattern element in order:
  \begin{itemize}
  \item If the pattern variable is already bound (from an earlier
    join step or a non-linear occurrence), look up the bound
    representative in the residual pool.  If found with sufficient
    multiplicity, subtract the required amount and continue.
    Otherwise, fail this branch.
  \item If unbound, iterate over all residual entries with nonzero
    multiplicity.  For each candidate, bind the variable to that
    representative, subtract the multiplicity, and recurse to the
    next pattern element.  On backtrack, restore the multiplicity.
  \end{itemize}

\item After all pattern elements are bound, handle the rest variable
  (if present): collect all residual entries with nonzero multiplicity
  into a multiset and bind the rest variable.  If no rest variable
  exists, verify that the residual is empty (exact match).
\end{enumerate}

\paragraph{Maximum partition semantics.}  When binding an unbound
variable to a residual entry, the engine takes the \emph{entire}
available multiplicity of that element, not just one copy.
This is the maximum partition: each distinct element is assigned to
exactly one pattern variable with its full multiplicity.  This avoids
enumerating the exponentially many ways to split copies of the same
element across variables.  The branching factor is the number of
distinct elements in the residual, not the total multiplicity.

Given term \texttt{(op a:4 b:3 c:2 d:1)} and pattern
\texttt{(op a:i y:j ..rest)}, where \texttt{a} is a ground e-class
and \texttt{y} is a pattern variable, the matcher:
\begin{enumerate}
\item Looks up ground \texttt{a} in the residual pool, binds
  \texttt{i}=4, removes $a^4$ from pool.
\item Enumerates unbound \texttt{y} over remaining elements:
  \begin{itemize}
  \item \texttt{y}=$b$, \texttt{j}=3 $\Rightarrow$
    \texttt{..rest}=$\{c^2, d^1\}$
  \item \texttt{y}=$c$, \texttt{j}=2 $\Rightarrow$
    \texttt{..rest}=$\{b^3, d^1\}$
  \item \texttt{y}=$d$, \texttt{j}=1 $\Rightarrow$
    \texttt{..rest}=$\{b^3, c^2\}$
  \end{itemize}
\end{enumerate}
The ground lookup is O(1); branching is only over the 3 remaining
distinct elements, independent of multiplicities.  For a multiset
with $d$ distinct elements and $p$ pattern variables, the search
space is $O(d^p)$, independent of multiplicities.

\paragraph{Multiplicity constraints.}  Each pattern element may
carry a constraint on its multiplicity (e.g., \texttt{a:k>=2}).
The constraint is checked when binding: if the available multiplicity
does not satisfy the constraint, the candidate is skipped.  If a
multiplicity variable appears multiple times in the pattern
(non-linear), all occurrences must bind to the same value.
Constraints compile to interval intersections; if the intersection
is empty, the entire query is statically unsatisfiable.

\paragraph{A and ACI variants.}  For A (associative, ordered)
operators, the decomposition enumerates contiguous subsequences
instead of sub-multisets: prefix and suffix rest variables capture
the unmatched head and tail.  For ACI (idempotent) operators,
multiplicities are always 1, so the decomposition simplifies to
subset enumeration.

\section{Proof Logging}

\subsection{Copy-on-First-Re-Canonization}

Rebuild re-canonizes nodes in place, destroying original child
pointers that proof chains need.  SP-Egraph resolves this with a
history bit on each e-node (the MSB of the operator field).  During
rebuild, before updating a node's children, the engine checks the
history bit.  If clear, the original children are saved to a proof
buffer and the bit is set.  If already set, the copy is skipped.
This ensures proof reconstruction can recover pre-merge structure at
a cost of at most one copy per node over its entire lifetime.

\subsection{Dual-Parent-Pointer Union-Find}

The union-find maintains two parent arrays.  The fast path uses path
compression for O($\alpha(n)$) find.  The proof path is never
compressed: it preserves the original merge tree.  Each union records
a justification (rewrite, congruence, or axiom) alongside the
proof-path edge.

\subsection{Euler-Tour LCA for Batch Extraction}

To explain $a \equiv b$, the engine finds their LCA in the proof
forest.  For single queries, a naive walk-up costs O(depth).  For
batch proof export, the Bender--Farach-Colton
algorithm~\cite{bfc} achieves O($n$) preprocessing and O(1) per
LCA query via Euler tour, exploiting the $\pm 1$ depth property for
block-decomposed RMQ.

Two implementations are provided: a naive walk-up for single queries
(O(depth)) and Euler-tour LCA for batch extraction (O($n$)
preprocessing, O(1) per query).

\subsection{Compile-Time Opt-Out}

A \texttt{const PROOFS: bool} generic on \texttt{EGraph} controls
all proof machinery.  When \texttt{false}, the compiler eliminates
proof arrays, history bits, justification vectors, and proof buffers
entirely.  The non-proof configuration has no residual overhead.

\section{Implementation and Status}

SP-Egraph is implemented in approximately 19{,}000 lines of Rust (plus 10{,}000 lines of tests).
The semi-persistent containers, e-graph core, pattern matching
engine, compilation pipeline, proof logging, and saturation loop are
fully implemented and tested.  The design is parameterized so that
every e-class reference (enode children, union-find links, match
bindings) can carry an edge label describing how bound variables
route between parent and child scopes.  The default label type is
unit (no binders, zero overhead).  Future work will instantiate this
with director matrices and slot maps to support alpha-equivalence,
following the approach of slotted e-graphs~\cite{slotted} and
director strings~\cite{sinot}.  Other planned extensions include
cost-based extraction via partial weighted Max-SAT and stratified
negation using the generational structure as stratum boundaries.

\section{Conclusion}

SP-Egraph's design rests on three contributions: pervasive
semi-persistence built from a single vector primitive with O(1)
snapshots and O($k$) restore; an extension of leapfrog triejoin to
A/AC/ACI matching with maximum partition semantics and multiplicity
constraints; and proof logging via copy-on-first-re-canonization
with Euler-tour LCA, with a compile-time generic that eliminates
all proof overhead when not needed.  Both semi-persistence and proof
logging are compile-time opt-out via \texttt{const} generics
(\texttt{TRACK} and \texttt{PROOFS}), with zero residual overhead
when disabled.  Semi-persistence is the
unifying mechanism: the same generational protocol that yields O(1)
snapshots also supplies stratum boundaries for stratified negation
and rollback for exploratory search.

\bibliographystyle{ACM-Reference-Format}
\begin{thebibliography}{10}

\bibitem{egg}
M.~Willsey, C.~Nandi, Y.~R. Wang, O.~Flatt, Z.~Tatlock, and
P.~Panchekha.
\newblock egg: Fast and extensible equality saturation.
\newblock In \emph{Proc.\ ACM Program.\ Lang.}, 5(POPL), 2021.

\bibitem{egglog}
Y.~Zhang, Y.~R. Wang, O.~Flatt, D.~Cao, P.~Zucker, E.~Roesner,
M.~Willsey, and Z.~Tatlock.
\newblock Better together: Unifying {D}atalog and equality saturation.
\newblock In \emph{Proc.\ ACM Program.\ Lang.}, 7(PLDI), 2023.

\bibitem{veldhuizen}
T.~L. Veldhuizen.
\newblock Leapfrog triejoin: A simple, worst-case optimal join
algorithm.
\newblock In \emph{Proc.\ ICDT}, 2014.

\bibitem{conchon-filliatre}
S.~Conchon and J.-C. Filli\^{a}tre.
\newblock Semi-persistent data structures.
\newblock In \emph{Proc.\ ESOP}, 2008.

\bibitem{iguernlala}
M.~Iguernlala.
\newblock \emph{Strengthening the Heart of an SMT-Solver: Design and
Implementation of Efficient Decision Procedures}.
\newblock PhD thesis, Universit\'{e} Paris-Sud, 2013.

\bibitem{bfc}
M.~A. Bender and M.~Farach-Colton.
\newblock The {LCA} problem revisited.
\newblock In \emph{Proc.\ LATIN}, 2000.

\bibitem{slotted}
R.~Schneider, M.~Rossel, A.~Shaikhha, A.~Goens, and M.~Steuwer.
\newblock Slotted e-graphs: First-class support for (bound) variables
in e-graphs.
\newblock In \emph{Proc.\ ACM Program.\ Lang.}, 9(PLDI), 2025.

\bibitem{sinot}
F.-R. Sinot, M.~Fern\'{a}ndez, and I.~Mackie.
\newblock Lambda-calculus with director strings.
\newblock \emph{Appl.\ Algebra Eng.\ Commun.\ Comput.}, 15:393--437,
2005.

\end{thebibliography}

\end{document}
